\chapter{Drop-Jump Validation Against Optoelectronics}
\label{ch:dropjump}

% SOURCE: results_dropjump_validation_v2.md
% STATUS: converted from markdown
% NOTE: [source:] tags stripped per ASSEMBLY_PLAN.md (harvested to appF)


This chapter presents ground-truth validation of the monocular camera pipeline (MediaPipe Heavy variant) on the OpenCap Drop-Jump dataset against synchronized 3D optoelectronic motion capture and force-plate ground truth.



\section{Temporal Synchronization and Event Detection}
\label{sec:6-1-temporal-synchronization-and-event-detection}


Joint trajectories from video and 3D motion capture were aligned across 48 trials (24 symmetric, 24 asymmetric landings across 8 subjects) using Ground Reaction Force (GRF)-anchored lag alignment.

\subsection{Synchronization Methodology}
\label{sec:6-1-1-synchronization-methodology}

The temporal anchor matched physical landing contact ($F_y > 20$ N) to kinematic flexion onset in video (local knee extension minimum). Native sample rates:
\begin{itemize}
  \item Synced video frame rate: \textbf{60.0000 FPS}.
  \item Optoelectronic Mocap IK sample rate: \textbf{100.0000 Hz}.
  \item Force plate sample rate: \textbf{2000.0000 Hz}.
\end{itemize}

\subsection{Biomarker Mapping}
\label{sec:6-1-2-biomarker-mapping}

Biomechanical events were identified via force plate and kinematic thresholds:
\begin{enumerate}
  \item \textbf{Initial Contact (IC1)}: First frame where vertical GRF exceeded $20$ N.
  \item \textbf{Peak Absorption (PA1)}: Frame containing maximum knee flexion between IC1 and takeoff (TO1).
  \item \textbf{Takeoff (TO1)}: First frame where vertical GRF dropped below $10$ N after IC1.
  \item \textbf{Final Landing Contact (IC2)}: First frame where vertical GRF exceeded $20$ N post-flight.
\end{enumerate}

These events mapped to primary knee flexion biomarkers:
\begin{itemize}
  \item \textbf{Contact Flexion (Biomarker \protect\#1)}: Knee flexion angle at \textbf{IC1}.
  \item \textbf{Peak Landing Flexion (Biomarker \protect\#2)}: Knee flexion angle at \textbf{PA1}.
  \item \textbf{Landing Range of Motion (ROM) (Biomarker \protect\#3)}: Excursion $\text{ROM} = \text{PA1} - \text{IC1}$.
  \item \textbf{Flexion Loading Rate (Biomarker \protect\#6)}: Average angular velocity during early landing ($IC1 \rightarrow \text{early absorption}$).
\end{itemize}



\section{Headline Finding: Deep-Flexion Constant Bias (Timing-Clean)}
\label{sec:6-2-headline-finding-deep-flexion-constant-bias-timing-clean}


Monocular measurement error does not scale monotonically with flexion depth across landing. Isolating static coordinate errors from dynamic synchronization lag, knee flexion was evaluated at peak absorption frames (PA1), where joint angular velocity is near zero ($\omega \approx 0$).

\subsection{Static-Peak Error Analysis}
\label{sec:6-2-1-static-peak-error-analysis}

Across $n = 96$ peak flexion points (48 trials $\times$ 2 limbs):
\begin{itemize}
  \item \textbf{Mean Deep-Flexion Bias (Video - IK)}: \textbf{$+10.52^\circ$}.
  \item \textbf{95\% Limits of Agreement (LoA)}: \textbf{$[-5.54^\circ, 26.58^\circ]$}.
\end{itemize}

\subsection{Error-vs-Depth Correlation}
\label{sec:6-2-2-error-vs-depth-correlation}

To test if overestimation is depth-dependent within active landing ($70^\circ\text{--}120^\circ$), correlations between 3D Mocap angle and error were evaluated:
\begin{itemize}
  \item Pearson correlation: $r = -0.1568$ ($p = 0.1271$).
  \item Spearman correlation: $\rho = -0.1905$ ($p = 0.0631$).
\end{itemize}

Because neither correlation was significant ($p > 0.05$), measurement error behaves as a \textbf{constant systematic positive bias} within the landing flexion band rather than a depth-dependent slope.

\subsection{Shallow Flexion Contrast}
\label{sec:6-2-3-shallow-flexion-contrast}

At initial contact (Biomarker \protect\#1, shallow flexion), the pipeline exhibits an underestimation bias of \textbf{$-6.69^\circ$}. This confirms systematic $+10.52^\circ$ overestimation is specific to deep flexion, forming a static, correctable offset.



\section{Per-Biomarker Agreement (Bland-Altman Analysis)}
\label{sec:6-3-per-biomarker-agreement-bland-altman-analysis}


Bland-Altman statistics (mean bias, 95\% LoA, Pearson $r$) across all $n = 48$ trials evaluate cohort agreement. Table 6.1 summarizes metrics across validatable biomarkers.


\begin{landscape}
\begin{table}[htbp]
\centering
\small
\caption{Bland-Altman Agreement Summary for Knee Flexion Biomarkers ($n = 48$ trials)}
\label{tab:bland-altman-summary}
\begin{tabular}{l c c c p{4.0cm} c p{4.5cm}}
\hline
\textbf{Biomarker} & \textbf{Video Mean} & \textbf{IK Mean} & \textbf{Bias (Video$-$IK)} & \textbf{95\% Limits of Agreement} & \textbf{Pearson $r$} & \textbf{Trustworthiness Verdict} \\
\hline
\textbf{\protect\#1 Contact Flexion} & $13.51^\circ$ & $20.20^\circ$ & $-6.69^\circ$ & $[-26.77^\circ,\ 13.39^\circ]$ & $0.3209$ & Accurate (low bias, moderate variance) \\
\textbf{\protect\#2 Peak Landing Flexion} & $120.23^\circ$ & $100.51^\circ$ & $+19.72^\circ$ & $[7.73^\circ,\ 31.71^\circ]$ & $0.8238$ & Biased-systematic (constant overestimation, low variance) \\
\textbf{\protect\#3 Landing ROM} & $106.72^\circ$ & $80.31^\circ$ & $+26.41^\circ$ & $[2.34^\circ,\ 50.48^\circ]$ & $0.4020$ & Biased-systematic (constant overestimation, high variance) \\
\textbf{\protect\#6 Flexion Loading Rate} & $286.14^\circ/\text{s}$ & $272.84^\circ/\text{s}$ & $+13.30^\circ/\text{s}$ & $[-115.92^\circ/\text{s},\ 142.51^\circ/\text{s}]$ & $0.6076$ & High-variance (moderate bias, high variance) \\
\textbf{\protect\#5 Asymmetry (IK-only)} & N/A & $2.07^\circ$ & N/A & N/A & N/A & \textbf{IK-only, not video-validated} (far-leg occlusion); mean$=2.07^\circ$ (SD$=2.06^\circ$) \\
\hline
\end{tabular}
\end{table}
\end{landscape}

Across biomarkers, contact flexion (\protect\#1) achieves high accuracy with minor underestimation ($-6.69^\circ$), whereas peak flexion (\protect\#2) and ROM (\protect\#3) carry systematic positive bias requiring calibration. Loading rate (\protect\#6) exhibits high variance from sampling jitter, and asymmetry (\protect\#5) is restricted to 3D IK reference data due to contralateral limb occlusion. See Appendix D for detailed biomarker-level interpretation.
Figure 6.1 presents the Bland-Altman agreement analysis for each biomarker.

\begin{figure}[htbp]
  \centering
  \includegraphics[width=0.95\textwidth]{fig6_1_bland_altman.png}
  \caption{Bland-Altman Agreement Plots for All Four Validatable Knee Flexion Biomarkers ($n=48$ trials). Each panel shows mean bias (solid line), 95\% Limits of Agreement (dashed), and individual trial scatter.}
  \label{fig:bland-altman-biomarkers}
\end{figure}

\begin{figure}[htbp]
  \centering
  \includegraphics[width=0.82\textwidth]{fig6_2_error_vs_depth.png}
  \caption{Headline Error vs.\ Flexion Depth. Scatter plot of video--IK measurement error as a function of peak flexion depth, demonstrating depth-independent systematic bias (Pearson $r=-0.16$, $p=0.13$).}
  \label{fig:error-vs-depth}
\end{figure}

\begin{figure}[htbp]
  \centering
  \includegraphics[width=0.72\textwidth]{fig6_3_peak_scatter.png}
  \caption{Static Peak Flexion Error Scatter. Per-trial systematic overestimation of peak flexion by monocular video relative to 3D IK, confirming the $+10.52^\circ$ timing-clean bias.}
  \label{fig:peak-scatter}
\end{figure}


\subsection{Peak-to-Peak vs. Timing-Clean Bias Reconciliation}
\label{sec:6-3-1-peak-to-peak-vs-timing-clean-bias-reconciliation}

A critical comparison emerges between the \textbf{$+10.52^\circ$ timing-clean peak bias} (Section 6.2) and the \textbf{$+19.72^\circ$ peak-to-peak bias} in Table 6.1:
\begin{itemize}
  \item \textit{Timing-Clean peak bias ($+10.52^\circ$)}: Frame-matched error at exact peak time ($t_{\text{peak}}$) of reference 3D motion capture.
  \item \textit{Peak-to-Peak bias ($+19.72^\circ$)}: Absolute maximum video trajectory value versus absolute maximum Mocap value, regardless of timing alignment.
\end{itemize}

Because 2D pose estimators experience frame-rate limits and dynamic tracking overshoot during high-velocity impact landings, unaligned peak-to-peak comparison captures temporal overshoot, inflating apparent bias by $+9.20^\circ$. For slower sagittal movements (squats, lunges) lacking impact velocities, this validation isolates the deep-flexion projection component ($\sigma^2_{\text{proj}}$), which the uncertainty-weighted framework (Chapter 8) transfers.



\section{Robustness to Movement Conditions}
\label{sec:6-4-robustness-to-movement-conditions}


To evaluate whether tracking accuracy varies with movement loading, static peak flexion errors were compared between symmetric and asymmetric landings.

\subsection{Symmetric vs. Asymmetric Landing Bias}
\label{sec:6-4-1-symmetric-vs-asymmetric-landing-bias}

\begin{itemize}
  \item \textbf{Symmetric Landings ($n = 48$ points)}: Mean peak overestimation bias \textbf{$+10.36^\circ$} (SD: $6.89^\circ$).
  \item \textbf{Asymmetric Landings ($n = 48$ points)}: Mean peak overestimation bias \textbf{$+10.68^\circ$} (SD: $9.39^\circ$).
\end{itemize}

\subsection{Methodological Implications}
\label{sec:6-4-2-methodological-implications}

The condition difference is negligible ($\Delta 0.32^\circ$). This near-identical error distribution demonstrates measurement bias is independent of kinetic landing asymmetry. Operating as a stable measurement property rather than a kinetic confound, this bias can be treated as a constant, correctable offset in downstream screening.



\section{Establishing the Bias Is Projection Error, Not Timing Artefact}
\label{sec:6-5-establishing-the-bias-is-projection-error-not-timing-artefact}


To confirm peak overestimation represents spatial projection distortion rather than temporal synchronization artefact or software modeling mismatch, three diagnostic tests were conducted.

\subsection{Dynamic Lag Test}
\label{sec:6-5-1-dynamic-lag-test}

We tested whether temporal trajectory misalignment caused deep-flexion error:
\begin{itemize}
  \item \textit{Frame-matched (GRF-aligned) MAE}: \textbf{$14.80^\circ$} (RMSE: $15.53^\circ$).
  \item \textit{Peak-matched MAE}: \textbf{$20.15^\circ$} (RMSE: $20.98^\circ$).
\end{itemize}

Forcing peak alignment worsened MAE by $5.35^\circ$ and RMSE by $5.45^\circ$. This rules out timing lag: peak-alignment forces video overshoot to match Mocap peak, increasing error and confirming spatial projection origin.

\subsection{Software Modeling Defect Test}
\label{sec:6-5-2-software-modeling-defect-test}

To rule out coordinate mismatches between clinical OpenSim models and superficial marker geometry, we compared:
\begin{enumerate}
  \item OpenSim joint coordinates from inverse kinematics output.
  \item Superficial 3-point marker trigonometric model (hip, knee, ankle) converted to clinical flexion.
\end{enumerate}

Mean offset between modeling definitions was \textbf{$1.64^\circ$} (RMSE: $21.79^\circ$ vs. $22.93^\circ$ against video), ruling out software definition mismatches as the source of $10.52^\circ$ overestimation.

\subsection{Sync Method Stability Analysis}
\label{sec:6-5-3-sync-method-stability-analysis}

Force-plate \textbf{GRF-anchored alignment} was compared against mathematical \textbf{RMSE-minimisation alignment}:
\begin{itemize}
  \item \textit{subject2 DJ1}: GRF lag = \textbf{3.00 ms (0 frames)} vs. RMSE lag = \textbf{-33.33 ms (-2 frames)}.
  \item \textit{subject2 DJAsym1}: GRF lag = \textbf{-241.17 ms (-14 frames)} vs. RMSE lag = \textbf{200.00 ms (12 frames)}.
  \item \textit{subject8 DJ1}: GRF lag = \textbf{-157.17 ms (-9 frames)} vs. RMSE lag = \textbf{33.33 ms (2 frames)}.
\end{itemize}

Because Mocap IK windows are trimmed ($\sim 1.0$ s), RMSE-minimization is unstable, choosing out-of-phase alignments (e.g., $+12$ frames, shifting video peak $0.4$ s before mocap peak). Conversely, GRF anchoring directly anchors physical impact, remaining stable across trials. Thus, GRF anchoring was adopted as standard, and pooled frame-level error-vs-depth correlation ($r = 0.3491$, $n = 3046$ frames) was demoted to a cautionary supplementary figure.
Figure 6.2 illustrates the naive pooled frame-level error plotted against knee flexion depth prior to timing diagnostic decomposition.

Diagnostics confirm deep-flexion overestimation stems from \textbf{sagittal-plane projection foreshortening}: as the knee flexes deeply outside the orthogonal camera plane, 2D projection systematically inflates apparent knee flexion angle.
Figure 6.3 demonstrates that when temporal artifacts are removed, static peak-flexion error is a constant scalar offset across flexion depths.



\section{Limitations}
\label{sec:6-6-limitations}


Limitations of dataset recording and camera configuration are documented:

\subsection{Dataset Recording Truncation (Stabilisation Time)}
\label{sec:6-6-1-dataset-recording-truncation-stabilisation-time}

Dynamic stabilisation (Biomarker \protect\#4, Time-to-Stabilisation) could not be computed. Video recordings truncate \textbf{$0.05\text{ s}$ to $0.2\text{ s}$} after final landing contact (IC2). Evaluating stabilisation requires a $0.5$ s quiet-stance window (30 frames at 60 FPS) with flexion standard deviation $< 1.5^\circ$; short recording lengths render this check unfeasible. This represents raw dataset truncation rather than pipeline failure.

\subsection{Contralateral Occlusion (Asymmetry)}
\label{sec:6-6-2-contralateral-occlusion-asymmetry}

Inter-limb asymmetry (Biomarker \protect\#5) cannot be resolved via monocular sagittal video. During deep landing absorption, the closer leg ($\sim 100\%$ tracking visibility) occludes the farther leg ($\sim 0\%$ visibility). As monocular tracking cannot resolve occluded joint landmarks, asymmetry is demoted to 3D Mocap reference data (Mocap mean: $2.07^\circ$, SD: $2.06^\circ$) and excluded from video decision rules.

\subsection{Time-Series Pooling Warnings}
\label{sec:6-6-3-time-series-pooling-warnings}

Pooled frame-level error-vs-depth displays high timing contamination and non-monotonic behavior ($r = 0.3491$). It serves as a supplementary warning against naive time-series pooling in single-camera validation.

\subsection{Kinematic Angle Conventions}
\label{sec:6-6-4-kinematic-angle-conventions}

This chapter evaluates knee flexion using \textbf{clinical flexion} ($0^\circ$ = extension, larger angle = deeper bend), contrasting with the \textbf{included angle} convention ($\approx 180^\circ$ = extension, smaller angle = deeper bend) in squat and lunge chapters. Raw angles are not directly comparable; the uncertainty-weighting framework (Chapter 8) resolves this by transferring validated error variances ($\sigma^2_{\text{proj}}$) rather than raw joint values.



\section{Figures and Provenance}
\label{sec:6-7-figures-and-provenance}


Publication-ready validation figures support this chapter:
\begin{itemize}
  \item \textbf{Figure 6.1: Bland-Altman agreement plots for peak knee flexion, contact flexion, ROM, and loading rate biomarkers}, comparing video-derived measurements against OpenCap motion-capture ground truth. The mean bias and 95\% Limits of Agreement are indicated by solid and dashed horizontal lines respectively.
  \item \textbf{Figure 6.2: Frame-level pipeline error plotted against knee flexion depth}, showing the naive pooled distribution that appears depth-dependent prior to timing diagnostic decomposition.
  \item \textbf{Figure 6.3: Static peak-flexion error plotted against ground-truth peak knee flexion depth.} The near-zero slope of the linear regression confirms that pipeline error is a constant scalar offset (approximately +10.5°) rather than a depth-dependent degradation.
\end{itemize}